\label{Chapter Experiment} 
\section{Experiment}
    \subsection{Final Dataset}
        In the final cycle, 10938 images were used for training and 2683 images for validating (\autoref{table:DatasetInformation}). 2464 images (with 80\% for training, 20\% for validating, and 0\% for testing) from all the sources listed above were used to generate 9856 augmented images. The original dataset has inconsistent number of labels for each class (\autoref{chart:Originaldatasetlabeldistribution}) but the splitting algorithm has done a great job to balance out the number labels per class for training set (\autoref{chart:Trainingsetlabeldistribution}) and validating set (\autoref{chart:Validatingsetlabeldistribution}). Additionally, there are 1301 background images (equivalent to 10\% of the dataset) chosen for each sub-set.

        For the number of instances of each class, there were 41 for training and around 9 for validating from the original dataset (\autoref{table:LabelDistribution} and \autoref{table:LabelDistribution1}). 
        After that, all of that expanded into 164 more labels per class for training and 36 for validating.
    \subsection{Final Hyperparameter Evolving}
        The configuration described in \autoref{table:Configuration} was used for the final evolution.
        The program took 10 days to finish running and return the results in \autoref{table:Hyperparameters}.
    \subsection{Final Training}
        For the last training session, the configuration listed in \autoref{table:Configuration} was employed. The YOLOv5x6 architecture and the customized version of hyperparameter are chosen with a preference for high accuracy.The results are shown in \autoref{fig:result1}, \autoref{fig:result2}, and \autoref{fig:result3}.
    \subsection{Deploy}
        For the sake of demonstration, \textbf{YOLOv5 Hololive Waifu Detection}\cite{Application} application was made. \textbf{Streamlit} is used for a basic GUI, and the application is hosted on the \textbf{Hugging Face Spaces} (\autoref{fig:WebApp}).
        Hugging Face Spaces offer a straightforward website to host machine learning demonstration applications directly. 
        Streamlit and Gradio are two fantastic Python SDKs that are supported by the platform. 
        Users may also construct static Spaces with JavaScript and HTML for more personalized demonstrations\cite{huggingfacespaces}.
    \end{multicols}
    \clearpage